\PaperTitle{烟幕干扰弹投放策略的时空协同优化}

针对无人机投放烟幕干扰弹以遮蔽圆柱形真目标的问题，本文建立导弹、无人机、烟幕弹与烟幕云团的统一三维运动学模型，并以“来袭导弹到真目标的全部视线均被烟幕阻断”作为有效遮蔽判据。区别于只考察目标中心视线的近似方法，本文将真目标视为具有实际尺寸的圆柱体：单云团情形利用烟幕球对导弹视点形成的切锥，将整圆柱遮蔽转化为上、下底圆周上的连续极值问题；多云团情形采用“对每条目标视线，至少存在一个有效云团与之相交”的联合覆盖口径，即 $\forall\text{-ray}\ \exists\text{-cloud}$，从而允许不同云团共同覆盖目标视域。有效遮蔽时长由连续时间区间的并集测度计算，重叠时段不重复计时，不连续时段可以累加。

对于问题一，在 $g=9.8\ \mathrm{m/s^2}$ 下求得投放点为 $(17620,0,1800)$ m，起爆点为 $(17188,0,1736.496)$ m，严格整圆柱遮蔽区间为 $[8.056445,\allowbreak{}9.448088]$ s，有效时长为 $1.391643$ s。目标中心视线近似给出 $1.435082$ s，比严格结果多计约 $3.121\%$。对于问题二，采用中心视线代理筛选、差分进化广域搜索、多起点无梯度局部精修和严格几何统一复算，得到当前最佳已知时长 $4.542880$ s。对于问题三，在同机三弹及投弹间隔约束下，得到联合遮蔽区间 $[2.764945,\allowbreak{}8.278596]$ s，总时长 $5.513651$ s；当前候选实质由前两枚弹完成，不能据此断言第三枚弹在全局最优策略中无用。对于问题四，三架无人机各投一弹形成三段接力遮蔽，联合总时长为 $10.138857$ s。对于问题五，以三枚导弹遮蔽时长之和为主目标，通过单弹候选库、分配组合、逐机块精修和边际填充构成 40 维分层求解，得到 $(D_1,\allowbreak{}D_2,\allowbreak{}D_3)=(10.250038,\allowbreak{}1.912852,\allowbreak{}1.741134)$ s，$J_{\mathrm{sum}}=13.904024$ s，公平性诊断量 $J_{\min}=1.741134$ s；最终五架无人机各使用一枚有效弹。各问题均进行了可行性、网格收敛、边界探针、独立视线—球相交和内核回归检查。由于问题二至问题五的优化均受评估预算限制，本文只将其结果称为经过严格复算的最佳已知策略，不作全局最优声明。


\noindent\textbf{关键词：}烟幕干扰；视线遮蔽；圆柱目标；联合覆盖；区间并集；非光滑优化
